% !PRINTED 345
\weberpage{345}

This expression, however, holds generally. For if $m_1$ and $m_2$ are
relatively prime and if
\[
m_1 = N(\mathfrak{m}_1), \quad m_2 = N(\mathfrak{m}_2),
\]
then
\[
m = m_1 m_2 = N(\mathfrak{m}_1 \mathfrak{m}_2).
\]

From a representation of $m_1$, $m_2$ as norms one can therefore derive
a representation of $m$ as a norm.

If conversely
\[
m_1 m_2 = N(\mathfrak{m}) = \mathfrak{m}\,\mathfrak{m}',
\]
where $\mathfrak{m}$, $\mathfrak{m}'$ are conjugate ideals, then, if $m_1$ and $m$ are divisible by
a prime ideal $\mathfrak{p}$, $m_1$ must also be divisible by $\mathfrak{p}\,\mathfrak{p}'$, and hence
$m_1 = N(\mathfrak{m}_1)$ as well, and likewise $m_2 = N(\mathfrak{m}_2)$.

Accordingly, if $m_1$ and $m_2$ are relatively prime:
\weq{7}{\psi(m_1)\,\psi(m_2) = \psi(m_1 m_2).}
On the other hand, if $n_1$ and $n_2$ run through the divisors of $m_1$ and $m_2$,
\[
\Sigma(\Delta,\, n_1)\, \Sigma(\Delta,\, n_2) = \Sigma(\Delta,\, n_1 n_2),
\]
and with this formula (6) is established in general. We
state this as a theorem:

4. \gesperrt{A natural number $m$ can be represented in the quadratic
field $\Omega$ with the fundamental number $\Delta$ in}
\[
\Sigma(\Delta,\, n)
\]
\gesperrt{different ways as the norm of an ideal,
when $n$ runs through all the divisors of $m$.}

\S\,94. \gesperrt{The quadratic law of reciprocity.}

From the decomposition of the primes in the quadratic field
\gesperrt{Dedekind} derives a new proof of the law of reciprocity
of quadratic residues\footnote{\gesperrt{Dirichlet-Dedekind}, Zahlentheorie, 4. Aufl. (S. 636, Anm.).}.

By Bd.\ II, \S\,167, 3. the necessary and sufficient
condition that the prime ideal $\mathfrak{p}$ dividing the prime $p$
(in any field $\Omega$) be of the first degree is this,
that for \gesperrt{every} integer $\omega$ of the field $\Omega$
\weq{1}{\omega^p \equiv \omega \ (\mathrm{mod}\ \mathfrak{p}),}
%CONT>>
